\uuid{nP4w}
\titre{EMV pour la loi de Laplace -- minimisation par la médiane empirique}
\niveau{L2}
\module{Probabilités et Statistique}
\chapitre{Statistique}
\sousChapitre{Estimation par maximum de vraisemblance}
\theme{Loi de Laplace, médiane empirique, minimisation de la somme des valeurs absolues}
\auteur{}
\datecreate{2026-03-30}
\organisation{AMSCC}
\difficulte{3}

\contenu{
	
	\texte{
		Soit $X_1, \ldots, X_n$ des variables aléatoires i.i.d. de densité commune
		\[
		f_\theta(x) = \frac{1}{2}e^{-|x-\theta|}, \quad \theta \in \mathbb{R}.
		\]
		On note $X_{(1)} \leq X_{(2)} \leq \cdots \leq X_{(n)}$ les statistiques d'ordre
		associées à l'échantillon.
	}
	
	\begin{enumerate}
		
		%------------------------------------------------------------
		\item
		\question{
			Écrire la vraisemblance $L(\theta)$ et la log-vraisemblance $\ell(\theta)$.
			Montrer que maximiser $\ell(\theta)$ est équivalent à minimiser
			\[
			g(\theta) = \sum_{i=1}^n |X_i - \theta|.
			\]
		}
		\reponse{
			La vraisemblance est
			\[
			L(\theta) = \prod_{i=1}^n \frac{1}{2} e^{-|X_i - \theta|}
			= \left(\frac{1}{2}\right)^{\!n} \exp\!\left(-\sum_{i=1}^n |X_i - \theta|\right).
			\]
			En prenant le logarithme :
			\[
			\ell(\theta) = -n\ln 2 - \sum_{i=1}^n |X_i - \theta| = -n\ln 2 - g(\theta).
			\]
			Puisque $-n\ln 2$ est une constante, maximiser $\ell(\theta)$ revient bien à
			minimiser $g(\theta)$.
		}
		
		%------------------------------------------------------------
		\item
		\question{
			Montrer que la fonction $\theta \mapsto |x - \theta|$ est convexe pour tout
			$x \in \mathbb{R}$, et en déduire que $g$ est convexe.
		}
		\reponse{
			Pour tout $x$ fixé, la fonction $h : \theta \mapsto |x - \theta|$ est l'opposé
			d'une valeur absolue translatée. Elle est convexe car, pour tous
			$\theta_1, \theta_2 \in \mathbb{R}$ et $\lambda \in [0,1]$ :
			\begin{align*}
				h(\lambda \theta_1 + (1-\lambda)\theta_2)
				&= |x - \lambda\theta_1 - (1-\lambda)\theta_2| \\
				&= |\lambda(x-\theta_1) + (1-\lambda)(x-\theta_2)| \\
				&\leq \lambda|x-\theta_1| + (1-\lambda)|x-\theta_2| \\
				&= \lambda h(\theta_1) + (1-\lambda)h(\theta_2).
			\end{align*}
			$g$ est une somme finie de fonctions convexes, donc $g$ est convexe.
			Tout minimum local de $g$ est donc un minimum global.
		}
		
		%------------------------------------------------------------
		\item
		\question{
			Montrer que $g$ est dérivable sur chaque intervalle ouvert
			$(X_{(k)}, X_{(k+1)})$ pour $k = 0, 1, \ldots, n$
			(avec la convention $X_{(0)} = -\infty$ et $X_{(n+1)} = +\infty$),
			et calculer $g'(\theta)$ sur chacun de ces intervalles en fonction de $k$.
		}
		\reponse{
			Sur l'intervalle $(X_{(k)}, X_{(k+1)})$, aucune observation ne coïncide avec
			$\theta$, donc chaque terme $|X_i - \theta|$ est différentiable en $\theta$.
			Plus précisément :
			\[
			\frac{\mathrm{d}}{\mathrm{d}\theta}|X_i - \theta|
			= \begin{cases} +1 & \text{si } X_i < \theta, \\ -1 & \text{si } X_i > \theta. \end{cases}
			\]
			Sur $(X_{(k)}, X_{(k+1)})$, exactement $k$ observations vérifient $X_i < \theta$
			et $n - k$ vérifient $X_i > \theta$, donc :
			\[
			g'(\theta) = k \cdot (+1) + (n-k) \cdot (-1) = k - (n-k) = 2k - n.
			\]
		}
		
		%------------------------------------------------------------
		\item
		\question{
			Étudier le signe de $g'(\theta)$ sur chaque intervalle et en déduire les
			variations de $g$. Montrer que $g$ est décroissante puis croissante, et
			identifier l'intervalle (ou le point) où $g$ atteint son minimum.
		}
		\reponse{
			Sur l'intervalle $(X_{(k)}, X_{(k+1)})$, on a $g'(\theta) = 2k - n$.
			\begin{itemize}
				\item Si $k < n/2$, alors $g'(\theta) < 0$ : $g$ est strictement décroissante.
				\item Si $k > n/2$, alors $g'(\theta) > 0$ : $g$ est strictement croissante.
				\item Si $k = n/2$ (possible seulement lorsque $n$ est pair), alors
				$g'(\theta) = 0$ : $g$ est constante sur $(X_{(n/2)}, X_{(n/2+1)})$.
			\end{itemize}
			Ainsi $g$ est décroissante jusqu'à l'indice $k = \lfloor n/2 \rfloor$, puis
			croissante à partir de $k = \lceil n/2 \rceil$. Le minimum est atteint sur
			l'intervalle (éventuellement réduit à un point) correspondant au passage
			$k = \lfloor n/2 \rfloor$ à $k = \lceil n/2 \rceil$.
		}
		
		%------------------------------------------------------------
		\item
		\question{
			Distinguer les deux cas $n$ impair et $n$ pair pour identifier précisément
			le ou les minimiseurs de $g$.
		}
		\reponse{
			\textbf{Cas $n = 2p+1$ impair.}
			
			On a $n/2 = p + 1/2$, donc $\lfloor n/2 \rfloor = p$ et $\lceil n/2 \rceil = p+1$.
			\begin{itemize}
				\item Sur $(X_{(p)}, X_{(p+1)})$ : $g'(\theta) = 2p - n = -1 < 0$ ($g$ décroissante).
				\item Sur $(X_{(p+1)}, X_{(p+2)})$ : $g'(\theta) = 2(p+1) - n = 1 > 0$ ($g$ croissante).
			\end{itemize}
			$g$ admet donc un minimum global unique en $\theta = X_{(p+1)}$.
			
			\textbf{Cas $n = 2p$ pair.}
			
			\begin{itemize}
				\item Sur $(X_{(p-1)}, X_{(p)})$ : $g'(\theta) = 2(p-1) - n = -2 < 0$ ($g$ décroissante).
				\item Sur $(X_{(p)}, X_{(p+1)})$ : $g'(\theta) = 2p - n = 0$ ($g$ constante).
				\item Sur $(X_{(p+1)}, X_{(p+2)})$ : $g'(\theta) = 2(p+1) - n = 2 > 0$ ($g$ croissante).
			\end{itemize}
			$g$ est minimale et constante sur le segment $[X_{(p)}, X_{(p+1)}]$ : tout point
			de cet intervalle est un minimiseur.
		}
		
		%------------------------------------------------------------
		\item
		\question{
			Rappeler la définition de la médiane empirique $\mathrm{med}(X_1, \ldots, X_n)$
			et conclure que l'estimateur du maximum de vraisemblance est
			\[
			\hat{\theta}_{MV} = \mathrm{med}(X_1, \ldots, X_n).
			\]
		}
		\reponse{
			La médiane empirique est définie par
			\[
			\mathrm{med}(X_1,\ldots,X_n) =
			\begin{cases}
				X_{(p+1)} & \text{si } n = 2p+1 \text{ (impair)}, \\[4pt]
				\dfrac{X_{(p)}+X_{(p+1)}}{2} & \text{si } n = 2p \text{ (pair)}.
			\end{cases}
			\]
			D'après les deux questions précédentes :
			\begin{itemize}
				\item Si $n$ est impair, l'unique minimiseur de $g$ est $X_{(p+1)} = \mathrm{med}(X_1,\ldots,X_n)$.
				\item Si $n$ est pair, tout point de $[X_{(p)}, X_{(p+1)}]$ minimise $g$,
				et la médiane empirique $\tfrac{1}{2}(X_{(p)}+X_{(p+1)})$ appartient à cet
				intervalle.
			\end{itemize}
			Dans les deux cas, la médiane empirique minimise $g(\theta) = \sum_{i=1}^n |X_i - \theta|$,
			ce qui maximise la log-vraisemblance. On conclut :
			\[
			\hat{\theta}_{MV} = \mathrm{med}(X_1, \ldots, X_n).
			\]
		}
		
	\end{enumerate}
}